\documentclass[10pt,a4paper]{article}
\usepackage[utf8]{inputenc}
\usepackage[top=2cm,bottom=2cm,left=2cm,right=2cm]{geometry}
\usepackage{amsmath}
\usepackage{amssymb}
\usepackage{array}
\usepackage[dvipsnames,svgnames]{xcolor}
\usepackage{hyperref}

% Anchor persamaan unik meski nomor di-reset per soal.
\newcounter{eqblock}
\newcommand{\reseteq}{%
  \stepcounter{eqblock}%
  \setcounter{equation}{0}%
  \renewcommand{\theHequation}{\arabic{eqblock}.\arabic{equation}}%
}

\newcommand{\dg}[1]{\textcolor{Red}{#1}}

\title{Tugas Minggu 12 : Invers dan Determinan}
\author{Matematika untuk Sistem Informasi}
\date{}

\begin{document}

\maketitle

\begin{enumerate}

%===========================================================
\item Diberikan matriks berikut:
\[
A = \begin{bmatrix} -2 & 3 \\ 2 & 1 \end{bmatrix}, \qquad
B = \begin{bmatrix} 1 & -3 & 0 \\ -2 & 4 & 1 \\ 5 & -2 & 2 \end{bmatrix}, \qquad
C = \begin{bmatrix} 1 & -2 & 3 & 1 \\ 5 & -9 & 6 & 3 \\ -1 & 2 & -6 & -2 \\ 2 & 8 & 6 & 1 \end{bmatrix}.
\]
Hitunglah nilai determinan matriks $A$, $B$, dan $C$ menggunakan metode ekspansi kofaktor (pilih salah satu kolom atau salah satu baris).
\begin{enumerate}

%-----------------------------------------------------------
\item $\det(A)$

\textbf{Jawab:}
\begin{itemize}
\item \textbf{Langkah 1: Ekspansi Sepanjang Baris Pertama}

Untuk matriks $2 \times 2$, minor setiap entri hanyalah satu bilangan: $M_{11} = 1$ dan $M_{12} = 2$. Dengan tanda $(+, -)$:
\reseteq
\begin{align}
\det(A) &= (-2) \cdot (+1) + 3 \cdot (-2) = -2 - 6 = -8
\end{align}
Hasil ini sama dengan rumus langsung $ad - bc = (-2)(1) - (3)(2) = -8$.
% final: \det(A) = -8
\end{itemize}

%-----------------------------------------------------------
\item $\det(B)$

\textbf{Jawab:}
\begin{itemize}
\item \textbf{Langkah 1: Pilih Baris Pertama}

Baris pertama $(1, -3, 0)$ memuat nol, sehingga hanya $M_{11}$ dan $M_{12}$ yang diperlukan:
\reseteq
\begin{align}
M_{11} &= \begin{vmatrix} 4 & 1 \\ -2 & 2 \end{vmatrix} = 4 \cdot 2 - 1 \cdot (-2) = 8 + 2 = 10 \\
M_{12} &= \begin{vmatrix} -2 & 1 \\ 5 & 2 \end{vmatrix} = (-2) \cdot 2 - 1 \cdot 5 = -4 - 5 = -9
\end{align}

\item \textbf{Langkah 2: Jumlahkan dengan Tanda $(+, -, +)$}

Kalikan setiap entri baris pertama dengan kofaktornya:
\begin{align}
\det(B) &= 1 \cdot (+M_{11}) + (-3) \cdot (-M_{12}) + 0 \cdot C_{13} \nonumber \\
&= 1(10) + (-3)(9) + 0 = 10 - 27 = -17
\end{align}
% final: \det(B) = -17
\end{itemize}

%-----------------------------------------------------------
\item $\det(C)$

\textbf{Jawab:}
\begin{itemize}
\item \textbf{Langkah 1: Tulis Ekspansi Baris Pertama}

Baris pertama $(1, -2, 3, 1)$ dengan tanda $(+, -, +, -)$ (lihat Contoh~2.7 Topik~10):
\reseteq
\begin{align}
\det(C) &= 1 \cdot M_{11} - (-2) \cdot M_{12} + 3 \cdot M_{13} - 1 \cdot M_{14}
\end{align}

\item \textbf{Langkah 2: Hitung Empat Minor $3 \times 3$}

Hapus baris 1 dan kolom yang bersesuaian, lalu ekspansi setiap minor sepanjang baris pertamanya:
\begin{align}
M_{11} &= \begin{vmatrix} -9 & 6 & 3 \\ 2 & -6 & -2 \\ 8 & 6 & 1 \end{vmatrix}
= -9(-6 + 12) - 6(2 + 16) + 3(12 + 48) = -54 - 108 + 180 = 18 \\
M_{12} &= \begin{vmatrix} 5 & 6 & 3 \\ -1 & -6 & -2 \\ 2 & 6 & 1 \end{vmatrix}
= 5(-6 + 12) - 6(-1 + 4) + 3(-6 + 12) = 30 - 18 + 18 = 30 \\
M_{13} &= \begin{vmatrix} 5 & -9 & 3 \\ -1 & 2 & -2 \\ 2 & 8 & 1 \end{vmatrix}
= 5(2 + 16) + 9(-1 + 4) + 3(-8 - 4) = 90 + 27 - 36 = 81 \\
M_{14} &= \begin{vmatrix} 5 & -9 & 6 \\ -1 & 2 & -6 \\ 2 & 8 & 6 \end{vmatrix}
= 5(12 + 48) + 9(-6 + 12) + 6(-8 - 4) = 300 + 54 - 72 = 282
\end{align}

\item \textbf{Langkah 3: Substitusi}

Masukkan keempat minor ke (1):
\begin{align}
\det(C) &= 1(18) + 2(30) + 3(81) - 1(282) = 18 + 60 + 243 - 282 = 39
\end{align}

\item \textbf{Langkah 4: Cek dengan OBE}

Operasi $R_2 - 5R_1$, $R_3 + R_1$, dan $R_4 - 2R_1$ tidak mengubah determinan. Setelah itu kolom pertama hanya memuat entri $1$, sehingga
\begin{align}
\det(C) &= \begin{vmatrix} 1 & -2 & 3 & 1 \\ 0 & 1 & -9 & -2 \\ 0 & 0 & -3 & -1 \\ 0 & 12 & 0 & -1 \end{vmatrix}
= 1 \cdot \begin{vmatrix} 1 & -9 & -2 \\ 0 & -3 & -1 \\ 12 & 0 & -1 \end{vmatrix}
= 1(3 - 0) + 12(9 - 6) = 3 + 36 = 39
\end{align}
(minor $3 \times 3$ diekspansi sepanjang kolom pertama). Hasilnya sama.
% final: \det(C) = 39
\end{itemize}

\end{enumerate}

%===========================================================
\item Tentukan semua nilai $a$ dan $b$ sehingga matriks $A$ dan $B$ berikut tidak memiliki invers.
\[
A = \begin{bmatrix} a + b - 1 & 0 \\ 0 & 3 \end{bmatrix}, \qquad
B = \begin{bmatrix} 5 & 0 \\ 0 & 2a - 3b - 7 \end{bmatrix}.
\]

\textbf{Jawab:}
\begin{itemize}
\item \textbf{Langkah 1: Syarat Tidak Memiliki Invers}

Matriks persegi tidak memiliki invers jika dan hanya jika determinannya nol (matriks singular). Kedua matriks adalah matriks diagonal, sehingga determinannya hasil kali entri diagonal:
\reseteq
\begin{align}
\det(A) &= 3(a + b - 1) = 0 \;\Rightarrow\; a + b = 1 \\
\det(B) &= 5(2a - 3b - 7) = 0 \;\Rightarrow\; 2a - 3b = 7
\end{align}
Jadi $A$ singular untuk setiap pasangan dengan $a + b = 1$, dan $B$ singular untuk setiap pasangan dengan $2a - 3b = 7$.

\item \textbf{Langkah 2: Eliminasi $a$}

Agar $A$ dan $B$ sekaligus tidak memiliki invers, (1) dan (2) harus dipenuhi bersamaan. Kalikan (1) dengan 2, lalu kurangkan (2):
\begin{align}
(2a + 2b) - (2a - 3b) &= 2 - 7 \nonumber \\
5b &= -5 \;\Rightarrow\; b = -1 \nonumber \\
a &= 1 - b = 2 \quad \text{(Dari (1))} \nonumber
\end{align}

\item \textbf{Langkah 3: Verifikasi}

Substitusi $a = 2$ dan $b = -1$:
\begin{align}
\det(A) &= 3(2 - 1 - 1) = 3 \cdot 0 = 0 \\
\det(B) &= 5(4 + 3 - 7) = 5 \cdot 0 = 0
\end{align}
Jadi $A$ dan $B$ sama-sama tidak memiliki invers tepat ketika $a = 2$ dan $b = -1$.
% final: a = 2
% final: b = -1
\end{itemize}

%===========================================================
\item Diketahui matriks berikut:
\[
A = \begin{bmatrix} 2 & 5 & 5 \\ -1 & -1 & 0 \\ 2 & 4 & 3 \end{bmatrix}, \qquad
B = \begin{bmatrix} \frac{1}{5} & \frac{1}{5} & -\frac{2}{5} \\[4pt] \frac{1}{5} & \frac{1}{5} & \frac{1}{10} \\[4pt] \frac{1}{5} & -\frac{4}{5} & \frac{1}{10} \end{bmatrix}.
\]
Tentukan invers matriks $A$ dan $B$ menggunakan matriks adjoin.
\begin{enumerate}

%-----------------------------------------------------------
\item $A^{-1}$

\textbf{Jawab:}
\begin{itemize}
\item \textbf{Langkah 1: Hitung Determinan}

Ekspansi sepanjang baris pertama $(2, 5, 5)$:
\reseteq
\begin{align}
\det(A) &= 2\begin{vmatrix} -1 & 0 \\ 4 & 3 \end{vmatrix} - 5\begin{vmatrix} -1 & 0 \\ 2 & 3 \end{vmatrix} + 5\begin{vmatrix} -1 & -1 \\ 2 & 4 \end{vmatrix}
= 2(-3) - 5(-3) + 5(-2) = -1
\end{align}
Karena $\det(A) = -1 \neq 0$, matriks $A$ memiliki invers.

\item \textbf{Langkah 2: Hitung Sembilan Kofaktor}

Gunakan pola tanda papan catur:
\begin{align}
C_{11} &= +\begin{vmatrix} -1 & 0 \\ 4 & 3 \end{vmatrix} = -3, &
C_{12} &= -\begin{vmatrix} -1 & 0 \\ 2 & 3 \end{vmatrix} = 3, &
C_{13} &= +\begin{vmatrix} -1 & -1 \\ 2 & 4 \end{vmatrix} = -2 \nonumber \\
C_{21} &= -\begin{vmatrix} 5 & 5 \\ 4 & 3 \end{vmatrix} = 5, &
C_{22} &= +\begin{vmatrix} 2 & 5 \\ 2 & 3 \end{vmatrix} = -4, &
C_{23} &= -\begin{vmatrix} 2 & 5 \\ 2 & 4 \end{vmatrix} = 2 \nonumber \\
C_{31} &= +\begin{vmatrix} 5 & 5 \\ -1 & 0 \end{vmatrix} = 5, &
C_{32} &= -\begin{vmatrix} 2 & 5 \\ -1 & 0 \end{vmatrix} = -5, &
C_{33} &= +\begin{vmatrix} 2 & 5 \\ -1 & -1 \end{vmatrix} = 3 \nonumber
\end{align}

\item \textbf{Langkah 3: Matriks Kofaktor dan Adjoin}

Susun kofaktor menurut letaknya, lalu transposkan:
\begin{align}
C &= \begin{bmatrix} -3 & 3 & -2 \\ 5 & -4 & 2 \\ 5 & -5 & 3 \end{bmatrix}
\quad\Rightarrow\quad
\operatorname{adj}(A) = C^T = \begin{bmatrix} -3 & 5 & 5 \\ 3 & -4 & -5 \\ -2 & 2 & 3 \end{bmatrix}
\end{align}

\item \textbf{Langkah 4: Bagi dengan Determinan}

Karena $\det(A) = -1$, invers adalah adjoin dengan semua tanda dibalik:
\begin{align}
A^{-1} &= \frac{1}{-1}\operatorname{adj}(A) = \begin{bmatrix} 3 & -5 & -5 \\ -3 & 4 & 5 \\ 2 & -2 & -3 \end{bmatrix}
\end{align}

\item \textbf{Langkah 5: Verifikasi}

Kalikan $A$ dengan hasil (3):
\begin{align}
AA^{-1} &= \begin{bmatrix}
6 - 15 + 10 & -10 + 20 - 10 & -10 + 25 - 15 \\
-3 + 3 + 0 & 5 - 4 + 0 & 5 - 5 + 0 \\
6 - 12 + 6 & -10 + 16 - 6 & -10 + 20 - 9
\end{bmatrix}
= \begin{bmatrix} 1 & 0 & 0 \\ 0 & 1 & 0 \\ 0 & 0 & 1 \end{bmatrix} = I
\end{align}
% final: A^{-1} = \begin{bmatrix} 3 & -5 & -5 \\ -3 & 4 & 5 \\ 2 & -2 & -3 \end{bmatrix}
\end{itemize}

%-----------------------------------------------------------
\item $B^{-1}$

\begingroup
\renewcommand{\arraystretch}{1.3}
\textbf{Jawab:}
\begin{itemize}
\item \textbf{Langkah 1: Hitung Determinan}

Ekspansi sepanjang baris pertama; minor yang diperlukan adalah
\reseteq
\begin{align}
M_{11} &= \begin{vmatrix} \frac{1}{5} & \frac{1}{10} \\ -\frac{4}{5} & \frac{1}{10} \end{vmatrix} = \frac{1}{50} + \frac{4}{50} = \frac{1}{10}, \qquad
M_{12} = \begin{vmatrix} \frac{1}{5} & \frac{1}{10} \\ \frac{1}{5} & \frac{1}{10} \end{vmatrix} = \frac{1}{50} - \frac{1}{50} = 0 \nonumber \\
M_{13} &= \begin{vmatrix} \frac{1}{5} & \frac{1}{5} \\ \frac{1}{5} & -\frac{4}{5} \end{vmatrix} = -\frac{4}{25} - \frac{1}{25} = -\frac{1}{5} \nonumber
\end{align}
sehingga
\begin{align}
\det(B) &= \frac{1}{5}\left(\frac{1}{10}\right) - \frac{1}{5}(0) + \left(-\frac{2}{5}\right)\left(-\frac{1}{5}\right) = \frac{1}{50} + \frac{4}{50} = \frac{1}{10} \neq 0
\end{align}
Cek dengan sifat $\det(pA) = p^n\det(A)$: tulis $B = \frac{1}{10}K$, maka
\[
K = \begin{bmatrix} 2 & 2 & -4 \\ 2 & 2 & 1 \\ 2 & -8 & 1 \end{bmatrix}, \qquad
\det(K) = 100, \qquad
\det(B) = \left(\frac{1}{10}\right)^3 \cdot 100 = \frac{1}{10}.
\]

\item \textbf{Langkah 2: Hitung Sembilan Kofaktor}

Baris pertama memakai minor dari Langkah 1; baris lainnya dihitung dengan cara yang sama:
\begin{align}
C_{11} &= +M_{11} = \frac{1}{10}, \qquad C_{12} = -M_{12} = 0, \qquad C_{13} = +M_{13} = -\frac{1}{5} \nonumber \\
C_{21} &= -\begin{vmatrix} \frac{1}{5} & -\frac{2}{5} \\ -\frac{4}{5} & \frac{1}{10} \end{vmatrix} = -\left(\frac{1}{50} - \frac{16}{50}\right) = \frac{3}{10} \nonumber \\
C_{22} &= +\begin{vmatrix} \frac{1}{5} & -\frac{2}{5} \\ \frac{1}{5} & \frac{1}{10} \end{vmatrix} = \frac{1}{50} + \frac{4}{50} = \frac{1}{10} \nonumber \\
C_{23} &= -\begin{vmatrix} \frac{1}{5} & \frac{1}{5} \\ \frac{1}{5} & -\frac{4}{5} \end{vmatrix} = -\left(-\frac{4}{25} - \frac{1}{25}\right) = \frac{1}{5} \nonumber \\
C_{31} &= +\begin{vmatrix} \frac{1}{5} & -\frac{2}{5} \\ \frac{1}{5} & \frac{1}{10} \end{vmatrix} = \frac{1}{50} + \frac{4}{50} = \frac{1}{10} \nonumber \\
C_{32} &= -\begin{vmatrix} \frac{1}{5} & -\frac{2}{5} \\ \frac{1}{5} & \frac{1}{10} \end{vmatrix} = -\frac{1}{10} \nonumber \\
C_{33} &= +\begin{vmatrix} \frac{1}{5} & \frac{1}{5} \\ \frac{1}{5} & \frac{1}{5} \end{vmatrix} = \frac{1}{25} - \frac{1}{25} = 0 \nonumber
\end{align}

\item \textbf{Langkah 3: Matriks Kofaktor dan Adjoin}

Susun kofaktor menurut letaknya, lalu transposkan:
\begin{align}
C &= \begin{bmatrix} \frac{1}{10} & 0 & -\frac{1}{5} \\ \frac{3}{10} & \frac{1}{10} & \frac{1}{5} \\ \frac{1}{10} & -\frac{1}{10} & 0 \end{bmatrix}
\quad\Rightarrow\quad
\operatorname{adj}(B) = C^T = \begin{bmatrix} \frac{1}{10} & \frac{3}{10} & \frac{1}{10} \\ 0 & \frac{1}{10} & -\frac{1}{10} \\ -\frac{1}{5} & \frac{1}{5} & 0 \end{bmatrix}
\end{align}

\item \textbf{Langkah 4: Bagi dengan Determinan}

Membagi dengan $\det(B) = \frac{1}{10}$ sama dengan mengalikan dengan $10$:
\begin{align}
B^{-1} &= \frac{1}{\frac{1}{10}}\operatorname{adj}(B) = 10\operatorname{adj}(B) = \begin{bmatrix} 1 & 3 & 1 \\ 0 & 1 & -1 \\ -2 & 2 & 0 \end{bmatrix}
\end{align}

\item \textbf{Langkah 5: Verifikasi}

Kalikan $B$ dengan hasil (3):
\begin{align}
BB^{-1} &= \begin{bmatrix}
\frac{1}{5} + 0 + \frac{4}{5} & \frac{3}{5} + \frac{1}{5} - \frac{4}{5} & \frac{1}{5} - \frac{1}{5} + 0 \\
\frac{1}{5} + 0 - \frac{1}{5} & \frac{3}{5} + \frac{1}{5} + \frac{1}{5} & \frac{1}{5} - \frac{1}{5} + 0 \\
\frac{1}{5} + 0 - \frac{1}{5} & \frac{3}{5} - \frac{4}{5} + \frac{1}{5} & \frac{1}{5} + \frac{4}{5} + 0
\end{bmatrix}
= \begin{bmatrix} 1 & 0 & 0 \\ 0 & 1 & 0 \\ 0 & 0 & 1 \end{bmatrix} = I
\end{align}
% final: B^{-1} = \begin{bmatrix} 1 & 3 & 1 \\ 0 & 1 & -1 \\ -2 & 2 & 0 \end{bmatrix}
\end{itemize}
\endgroup

\end{enumerate}

\end{enumerate}

\end{document}
